%% toc: P2 17 17.4
\section{Mendelson's first-order theories and the deduction theorem}
\label{sec:p02-w02c-p02-017-mendelson}

\emph{Purpose.}  This section presents, in Hilbert style, the first-order
theories of Mendelson's \emph{Introduction to Mathematical Logic}, fourth edition,
chapter~2, in the source's own notation: the logical axiom schemas (A1)--(A5)
with the side conditions of (A4) and (A5), the rules modus ponens and
generalization, the dependence of a step of a deduction on a hypothesis, and
the deduction theorem with Mendelson's hypothesis on applications of
generalization \cite[\S\S 2.3--2.4, pp.~69--75]{Mendelson1997}.  It is a
separate system from the registered first-order Hilbert calculus of
\cref{sec:p02-hilbert-fol} and from Kleene's predicate calculus of
\cref{sec:p02-w02-p02-025-kleene}; the differences are stated in
\cref{sec:p02-w02c-p02-017-relations}.

\emph{Inputs and conclusions.}  The input is a first-order language
\(\mathscr L\), a first-order theory \(K\) in \(\mathscr L\) (its proper
axioms), a set \(\Gamma\) of wfs of \(K\), a wf \(\mathscr B\), and a
deduction of a wf \(\mathscr C\) from \(\Gamma\) together with
\(\mathscr B\), given with a justification for each step.  The conclusions
are: \(\Gamma\vdash_K\mathscr C\) when \(\mathscr C\) does not depend upon
\(\mathscr B\) (\cref{thm:p02-w02c-p02-017-prop24}); and
\(\Gamma\vdash_K\mathscr B\Rightarrow\mathscr C\) when no application of
generalization to a wf depending upon \(\mathscr B\) quantifies a variable
free in \(\mathscr B\) (\cref{thm:p02-w02c-p02-017-prop25}), with two
corollaries that check the hypothesis more coarsely
(\cref{thm:p02-w02c-p02-017-cor26,thm:p02-w02c-p02-017-cor27}).  Without
such a hypothesis the deduction theorem fails
(\cref{thm:p02-w02c-p02-017-failure}).

\emph{Notation.}  The notation is Mendelson's: \(\neg\) and
\(\Rightarrow\) are the primitive connectives, \((\forall x_i)\) is the
universal quantifier, \(x_1,x_2,\ldots\) are the individual variables,
\(A^n_j\), \(f^n_j\) and \(a_j\) are predicate letters, function letters and
individual constants, script capitals \(\mathscr B,\mathscr C,\mathscr D,
\mathscr E\) range over wfs, and \(\Gamma\vdash_K\mathscr C\) says that
\(\mathscr C\) is a consequence of \(\Gamma\) in \(K\).  In this section
\(\vdash_K\) is Mendelson's relation, not the derivability relation of the
registered sheet.  Page numbers refer to the printed fourth edition.

\subsection{Languages, terms and the free-for relation}
\label{sec:p02-w02c-p02-017-language}

\begin{hbdefinition}{def:p02-w02c-p02-017-language}{Languages, terms and wfs in Mendelson}
A \emph{first-order language} \(\mathscr L\) has as symbols the connectives
\(\neg\) and \(\Rightarrow\), the quantifier symbol \(\forall\), the
parentheses and the comma, denumerably many individual variables
\(x_1,x_2,\ldots\), a finite or denumerable (possibly empty) set of function
letters \(f^n_j\), a finite or denumerable (possibly empty) set of individual
constants \(a_j\), and a nonempty set of predicate letters \(A^n_j\)
\cite[\S 2.2, p.~57]{Mendelson1997}.  The superscript \(n\) is the number of
arguments.
\begin{enumerate}
  \item \emph{Terms}: variables and individual constants are terms; if
        \(f^n_j\) is a function letter and \(t_1,\ldots,t_n\) are terms, then
        \(f^n_j(t_1,\ldots,t_n)\) is a term; nothing else is a term
        \cite[\S 2.1, p.~52]{Mendelson1997}.
  \item \emph{Wfs}: if \(A^n_j\) is a predicate letter and
        \(t_1,\ldots,t_n\) are terms, \(A^n_j(t_1,\ldots,t_n)\) is an
        (atomic) wf; if \(\mathscr B\) and \(\mathscr C\) are wfs and \(y\) is
        a variable, then \((\neg\mathscr B)\), \((\mathscr B\Rightarrow
        \mathscr C)\) and \(((\forall y)\mathscr B)\) are wfs; nothing else is
        a wf.  In \(((\forall y)\mathscr B)\) the wf \(\mathscr B\) is the
        \emph{scope} of the quantifier \((\forall y)\)
        \cite[\S 2.1, p.~52]{Mendelson1997}.  The other connectives and
        \((\exists x)\mathscr B\), which abbreviates
        \(\neg(\forall x)\neg\mathscr B\), are defined, not primitive
        \cite[\S 2.1, p.~52]{Mendelson1997}.
  \item An occurrence of a variable \(x\) in a wf \(\mathscr B\) is
        \emph{bound} if it is the occurrence in a quantifier
        \((\forall x)\) of \(\mathscr B\) or lies within the scope of a
        quantifier \((\forall x)\) of \(\mathscr B\); otherwise it is
        \emph{free}.  A variable is free (bound) in \(\mathscr B\) if it has a
        free (bound) occurrence there \cite[\S 2.1, p.~53]{Mendelson1997}.
        A wf with no free variables is \emph{closed}
        \cite[\S 2.2, p.~58]{Mendelson1997}.
  \item Writing \(\mathscr B\) as \(\mathscr B(x_i)\) indicates that
        \(x_i\) may be free in \(\mathscr B\) (it need not be, and other
        variables may be free as well); then \(\mathscr B(t)\) is the result
        of putting the term \(t\) for every free occurrence, if any, of
        \(x_i\) in \(\mathscr B\) \cite[\S 2.1, p.~54]{Mendelson1997}.  This
        is replacement of occurrences, with no renaming of bound variables.
  \item A term \(t\) is \emph{free for \(x_i\) in \(\mathscr B\)} if no free
        occurrence of \(x_i\) in \(\mathscr B\) lies within the scope of a
        quantifier \((\forall x_j)\) of \(\mathscr B\) for which \(x_j\) is a
        variable of \(t\) \cite[\S 2.1, p.~54]{Mendelson1997}.  Then no
        variable of \(t\) becomes bound when \(\mathscr B(t)\) is formed.  In
        particular \(x_i\) is free for \(x_i\) in every wf, a term without
        variables is free for every variable in every wf, and every term is
        free for \(x_i\) in a wf in which \(x_i\) has no free occurrence
        \cite[\S 2.1, p.~54]{Mendelson1997}.
\end{enumerate}
\end{hbdefinition}

\subsection{First-order theories and deductions}
\label{sec:p02-w02c-p02-017-theories}

\begin{hbdefinition}{def:p02-w02c-p02-017-theory}{First-order theory (Mendelson)}
Let \(\mathscr L\) be a first-order language.  A \emph{first-order theory}
\(K\) in \(\mathscr L\) has the symbols and wfs of \(\mathscr L\), and its
axioms and rules are these \cite[\S 2.3, pp.~69--70]{Mendelson1997}.

\emph{Logical axioms.}  For all wfs \(\mathscr B,\mathscr C,\mathscr D\) of
\(\mathscr L\):
\[
\begin{array}{ll}
 \text{(A1)} & \mathscr B\Rightarrow(\mathscr C\Rightarrow\mathscr B)\\
 \text{(A2)} & (\mathscr B\Rightarrow(\mathscr C\Rightarrow\mathscr D))
               \Rightarrow((\mathscr B\Rightarrow\mathscr C)
               \Rightarrow(\mathscr B\Rightarrow\mathscr D))\\
 \text{(A3)} & (\neg\mathscr C\Rightarrow\neg\mathscr B)
               \Rightarrow((\neg\mathscr C\Rightarrow\mathscr B)
               \Rightarrow\mathscr C)\\
 \text{(A4)} & (\forall x_i)\mathscr B(x_i)\Rightarrow\mathscr B(t)\\
 \text{(A5)} & (\forall x_i)(\mathscr B\Rightarrow\mathscr C)
               \Rightarrow(\mathscr B\Rightarrow(\forall x_i)\mathscr C)
\end{array}
\]
Side conditions: in (A4), \(\mathscr B(x_i)\) is a wf of \(\mathscr L\) and
\(t\) is a term of \(\mathscr L\) that is free for \(x_i\) in
\(\mathscr B(x_i)\); in (A5), \(\mathscr B\) contains no free occurrence of
\(x_i\).  A wf of either shape that fails its side condition is not an
axiom.  Since \(x_i\) is free for \(x_i\) in every wf, every wf
\((\forall x_i)\mathscr B\Rightarrow\mathscr B\) is an instance of (A4)
\cite[\S 2.3, p.~69]{Mendelson1997}.

\emph{Proper axioms.}  These vary from theory to theory.  A first-order
theory with no proper axioms is a \emph{first-order predicate calculus}
\cite[\S 2.3, p.~70]{Mendelson1997}.

\emph{Rules of inference.}
\begin{enumerate}
  \item \emph{Modus ponens} (MP): \(\mathscr C\) follows from \(\mathscr B\)
        and \(\mathscr B\Rightarrow\mathscr C\).
  \item \emph{Generalization} (Gen): \((\forall x_i)\mathscr B\) follows
        from \(\mathscr B\), for every variable \(x_i\)
        \cite[\S 2.3, p.~70]{Mendelson1997}.
\end{enumerate}
Gen carries no side condition.
\end{hbdefinition}

\begin{hbdefinition}{def:p02-w02c-p02-017-deduction}{Deduction and consequence}
The proof-theoretic vocabulary is the one Mendelson fixes for formal theories
in general and carries over to first-order theories
\cite[\S 1.4, pp.~34--35; \S 2.3, p.~69]{Mendelson1997}.
\begin{enumerate}
  \item A \emph{deduction} (proof) of \(\mathscr C\) from a set \(\Gamma\) of
        wfs in \(K\) is a finite sequence \(\mathscr D_1,\ldots,\mathscr D_n\)
        of wfs with \(\mathscr D_n=\mathscr C\) in which each
        \(\mathscr D_i\) is an axiom of \(K\), or belongs to \(\Gamma\), or
        follows from preceding members by MP or Gen.  The members of
        \(\Gamma\) are the hypotheses (premisses).  \(\Gamma\vdash_K\mathscr C\)
        says that such a deduction exists, and \(\mathscr C\) is then a
        \emph{consequence} of \(\Gamma\) in \(K\); \(\vdash_K\mathscr C\)
        abbreviates \(\varnothing\vdash_K\mathscr C\), and then
        \(\mathscr C\) is a \emph{theorem} of \(K\)
        \cite[\S 1.4, pp.~34--35]{Mendelson1997}.  The subscript \(K\) is
        dropped when the theory is fixed \cite[\S 2.4, p.~71]{Mendelson1997},
        and \(\Gamma,\mathscr B\) stands for \(\Gamma\cup\{\mathscr B\}\).
  \item A \emph{justification} of a deduction from \(\Gamma\) assigns to each
        step one reason: Hyp (the wf is the member of \(\Gamma\) named by the
        step), an axiom of \(K\), MP from named earlier steps
        \(\mathscr D_j\) and \(\mathscr D_k=(\mathscr D_j\Rightarrow
        \mathscr D_i)\), or Gen from a named earlier step \(\mathscr D_j\)
        with a named quantified variable \(x_m\), so that
        \(\mathscr D_i=(\forall x_m)\mathscr D_j\).  The same sequence of wfs
        may admit several justifications; the notions below are relative to
        the one given.
\end{enumerate}
Because Gen is unrestricted, \(\Gamma\vdash_K\mathscr C\) allows
generalization on a variable that is free in a hypothesis: for every wf
\(\mathscr B\) and variable \(x_i\), the two-step sequence \(\mathscr B\)
(Hyp), \((\forall x_i)\mathscr B\) (Gen) shows
\(\mathscr B\vdash_K(\forall x_i)\mathscr B\)
\cite[\S 2.4, p.~73]{Mendelson1997}.  Three facts are used throughout: an
axiom or a member of \(\Gamma\) is a consequence of \(\Gamma\) by a
one-step deduction; if \(\Gamma\vdash_K\mathscr B\) and
\(\Gamma\vdash_K\mathscr B\Rightarrow\mathscr C\), then
\(\Gamma\vdash_K\mathscr C\), by concatenating the two deductions and
appending \(\mathscr C\) by MP; and if \(\Gamma\vdash_K\mathscr B\), then
\(\Gamma\vdash_K(\forall x_i)\mathscr B\), by appending one Gen step.
\end{hbdefinition}

\begin{hbwarning}{def:p02-w02c-p02-017-provisos}{The side conditions of (A4) and (A5)}
Both side conditions are needed for the axioms to be logically valid
\cite[\S 2.3, pp.~70--71]{Mendelson1997}.  The two pseudo-instances below
are Mendelson's; the evaluations use the structures and satisfaction
relation of \cref{sec:fol-semantics-01,sec:fol-semantics-02}, reading
\(\Rightarrow\) as \(\to\) and each predicate letter as a relation symbol.
\begin{enumerate}
  \item \emph{(A4) without ``free for''.}  Let \(\mathscr B(x_1)\) be
        \(\neg(\forall x_2)A^2_1(x_1,x_2)\) and \(t\) be \(x_2\); the free
        occurrence of \(x_1\) lies in the scope of \((\forall x_2)\), so
        \(x_2\) is not free for \(x_1\).  The wf
        \((\forall x_1)\neg(\forall x_2)A^2_1(x_1,x_2)\Rightarrow
        \neg(\forall x_2)A^2_1(x_2,x_2)\) has the shape of (A4) but is not an
        instance.  In a structure with at least two elements that interprets
        \(A^2_1\) as identity, every element differs from some element, so
        the antecedent is satisfied by every assignment, while
        \((\forall x_2)A^2_1(x_2,x_2)\) is satisfied too, so the consequent
        fails under every assignment.
  \item \emph{(A5) with \(x_i\) free in \(\mathscr B\).}  Let both
        \(\mathscr B\) and \(\mathscr C\) be \(A^1_1(x_1)\).  The wf
        \((\forall x_1)(A^1_1(x_1)\Rightarrow A^1_1(x_1))\Rightarrow
        (A^1_1(x_1)\Rightarrow(\forall x_1)A^1_1(x_1))\) has the shape of
        (A5) but is not an instance.  Over the integers with \(A^1_1\)
        interpreted as evenness, an assignment sending \(x_1\) to an even
        number satisfies the antecedent and \(A^1_1(x_1)\) but not
        \((\forall x_1)A^1_1(x_1)\).
\end{enumerate}
Neither wf is therefore true in every structure under every assignment, while
every theorem of a predicate calculus is
(\cref{thm:p02-w02c-p02-017-soundness}); admitting either as an axiom would
destroy that property.
\end{hbwarning}

\subsection{Dependence and the deduction theorem}
\label{sec:p02-w02c-p02-017-deduction}

The proof of the deduction theorem needs \(\vdash_K\mathscr B\Rightarrow
\mathscr B\).  Mendelson's proof obtains it from his
Proposition~2.1 (instances of tautologies) \cite[\S 2.4, p.~74]{Mendelson1997};
here it is derived directly from (A1), (A2) and MP, as in his Lemma~1.8 for
the propositional system \(L\) \cite[\S 1.4, p.~36]{Mendelson1997}, whose
axioms (A1)--(A3) are those of \(K\) \cite[\S 1.4, p.~35]{Mendelson1997}.

\begin{hblemma}{thm:p02-w02c-p02-017-identity}{proved=annex:p02-w02c-p02-017-identity}{Identity in a first-order theory}
For every first-order theory \(K\) and every wf \(\mathscr B\) of \(K\),
\(\vdash_K\mathscr B\Rightarrow\mathscr B\) by a five-step deduction that
uses only instances of (A1) and (A2) and MP, and no hypothesis and no Gen.
\end{hblemma}

\begin{hbdefinition}{def:p02-w02c-p02-017-dependence}{Dependence upon a hypothesis}
Let \(\mathscr D_1,\ldots,\mathscr D_n\) be a deduction from a set \(\Delta\)
of wfs, with a justification for each step, and let \(\mathscr B\in\Delta\).
The step \(\mathscr D_i\) \emph{depends upon} \(\mathscr B\) in this deduction
if and only if \cite[\S 2.4, p.~73]{Mendelson1997}:
\begin{enumerate}
  \item \(\mathscr D_i\) is \(\mathscr B\) and its justification is Hyp; or
  \item \(\mathscr D_i\) is justified by MP or Gen from preceding steps at
        least one of which depends upon \(\mathscr B\).
\end{enumerate}
The definition is by recursion on \(i\), and whether \(\mathscr D_i\)
depends upon \(\mathscr B\) is fixed by \(\mathscr D_1,\ldots,\mathscr D_i\)
and their justifications.  So an initial segment
\(\mathscr D_1,\ldots,\mathscr D_j\) is itself a justified deduction in
which each of its steps has the same dependences as in the whole.
\end{hbdefinition}

\begin{hbtheorem}{thm:p02-w02c-p02-017-prop24}{proved=annex:p02-w02c-p02-017-prop24}{Mendelson's Proposition 2.4}
If \(\mathscr C\) does not depend upon \(\mathscr B\) in a deduction showing
that \(\Gamma,\mathscr B\vdash_K\mathscr C\), then
\(\Gamma\vdash_K\mathscr C\) \cite[\S 2.4, p.~74]{Mendelson1997}.
\end{hbtheorem}

\begin{hbtheorem}{thm:p02-w02c-p02-017-prop25}{proved=annex:p02-w02c-p02-017-prop25}{Mendelson's Proposition 2.5, the deduction theorem}
Suppose that, in some deduction showing that
\(\Gamma,\mathscr B\vdash_K\mathscr C\), no application of Gen to a wf that
depends upon \(\mathscr B\) has as its quantified variable a free variable of
\(\mathscr B\).  Then \(\Gamma\vdash_K\mathscr B\Rightarrow\mathscr C\)
\cite[\S 2.4, p.~74]{Mendelson1997}.
\end{hbtheorem}

The hypothesis concerns only those Gen steps whose premiss depends upon
\(\mathscr B\); a Gen step on a variable free in \(\mathscr B\) is allowed
when its premiss does not depend upon \(\mathscr B\).  Mendelson finds the
hypothesis unwieldy in practice and follows it with two corollaries whose
hypotheses are easier to check and imply it
\cite[\S 2.4, pp.~74--75]{Mendelson1997}.

\begin{hbcorollary}{thm:p02-w02c-p02-017-cor26}{proved=annex:p02-w02c-p02-017-cor26}{Mendelson's Corollary 2.6}
If a deduction showing that \(\Gamma,\mathscr B\vdash_K\mathscr C\) contains
no application of Gen whose quantified variable is free in \(\mathscr B\),
then \(\Gamma\vdash_K\mathscr B\Rightarrow\mathscr C\)
\cite[\S 2.4, p.~75]{Mendelson1997}.
\end{hbcorollary}

\begin{hbcorollary}{thm:p02-w02c-p02-017-cor27}{proved=annex:p02-w02c-p02-017-cor27}{Mendelson's Corollary 2.7}
If \(\mathscr B\) is a closed wf and \(\Gamma,\mathscr B\vdash_K\mathscr C\),
then \(\Gamma\vdash_K\mathscr B\Rightarrow\mathscr C\)
\cite[\S 2.4, p.~75]{Mendelson1997}.
\end{hbcorollary}

To use the deduction theorem twice on one deduction, one must know which Gen
steps the first application produces.  Mendelson states this as a
supplement that he counts as part of Propositions 2.4--2.7 and leaves the
proof as Exercise~2.29 \cite[\S 2.4, pp.~75--76]{Mendelson1997}.  It is
stated here existentially; the deductions constructed in the annex proofs of
the four results are the witnesses.

\begin{hbtheorem}{thm:p02-w02c-p02-017-extension}{proved=annex:p02-w02c-p02-017-extension}{Extension of Propositions 2.4--2.7}
Let a justified deduction showing that \(\Gamma,\mathscr B\vdash_K\mathscr C\)
be given, satisfying the hypothesis of
\cref{thm:p02-w02c-p02-017-prop24}, \cref{thm:p02-w02c-p02-017-prop25},
\cref{thm:p02-w02c-p02-017-cor26} or \cref{thm:p02-w02c-p02-017-cor27}.
Then there is a justified deduction of the corresponding conclusion
(\(\mathscr C\) for Proposition 2.4, \(\mathscr B\Rightarrow\mathscr C\) for
the other three) from \(\Gamma\) with this property: whenever it applies Gen
to a wf that depends upon a member \(\mathscr E\) of \(\Gamma\), the given
deduction contains an application of Gen with the same quantified variable to
a wf that depends upon \(\mathscr E\) \cite[\S 2.4, p.~75]{Mendelson1997}.
\end{hbtheorem}

\begin{hbexample}{def:p02-w02c-p02-017-example}{Distributing the universal quantifier over \(\Rightarrow\)}
Let \(\mathscr B,\mathscr C\) be wfs, and put
\(\mathscr H_1=(\forall x_1)(\mathscr B\Rightarrow\mathscr C)\) and
\(\mathscr H_2=(\forall x_1)\mathscr B\).  The following justified deduction
shows \(\mathscr H_1,\mathscr H_2\vdash_K(\forall x_1)\mathscr C\) in every
first-order theory \(K\).  The last column lists the hypotheses each step
depends upon.
\[
\begin{array}{rlll}
 1. & (\forall x_1)(\mathscr B\Rightarrow\mathscr C) & \text{Hyp} & \mathscr H_1\\
 2. & (\forall x_1)\mathscr B & \text{Hyp} & \mathscr H_2\\
 3. & (\forall x_1)(\mathscr B\Rightarrow\mathscr C)\Rightarrow
      (\mathscr B\Rightarrow\mathscr C) & \text{(A4)} & \text{none}\\
 4. & \mathscr B\Rightarrow\mathscr C & 1,3,\ \text{MP} & \mathscr H_1\\
 5. & (\forall x_1)\mathscr B\Rightarrow\mathscr B & \text{(A4)} & \text{none}\\
 6. & \mathscr B & 2,5,\ \text{MP} & \mathscr H_2\\
 7. & \mathscr C & 4,6,\ \text{MP} & \mathscr H_1,\mathscr H_2\\
 8. & (\forall x_1)\mathscr C & 7,\ \text{Gen} & \mathscr H_1,\mathscr H_2
\end{array}
\]
Steps 3 and 5 are instances of (A4) with \(t\) the variable \(x_1\) itself.
The only Gen step is step~8, with quantified variable \(x_1\), and \(x_1\) is
not free in \(\mathscr H_2\), where every occurrence of \(x_1\) is bound.  By
\cref{thm:p02-w02c-p02-017-cor26} with \(\Gamma=\{\mathscr H_1\}\),
\(\mathscr H_1\vdash_K\mathscr H_2\Rightarrow(\forall x_1)\mathscr C\).  By
\cref{thm:p02-w02c-p02-017-extension} there is a deduction of
\(\mathscr H_2\Rightarrow(\forall x_1)\mathscr C\) from \(\{\mathscr H_1\}\)
in which every Gen step applied to a wf depending upon \(\mathscr H_1\) has
quantified variable \(x_1\), because step~8 is the only Gen step of the
deduction above.  Since \(x_1\) is not free in \(\mathscr H_1\), that deduction
satisfies the hypothesis of \cref{thm:p02-w02c-p02-017-prop25} with
\(\Gamma=\varnothing\), and so
\[
 \vdash_K(\forall x_1)(\mathscr B\Rightarrow\mathscr C)\Rightarrow
 ((\forall x_1)\mathscr B\Rightarrow(\forall x_1)\mathscr C).
\]
This is Mendelson's Exercise~2.27(a) for the variable \(x_1\)
\cite[\S 2.4, p.~75]{Mendelson1997}.  Corollary~2.7 does not give the second
step directly when \(\mathscr H_1\) has free variables.
\end{hbexample}

\subsection{Failure without the hypothesis on Gen}
\label{sec:p02-w02c-p02-017-failure}

The failure is shown semantically, with the structures, assignments and
satisfaction relation of \cref{sec:fol-semantics-01,sec:fol-semantics-02}.
A wf of \(\mathscr L\) is read as a formula of the signature whose relation,
function and constant symbols are the predicate letters, function letters and
individual constants of \(\mathscr L\), with \(\Rightarrow\) read as \(\to\).
Mendelson's own semantics uses denumerable sequences in place of assignments
\cite[\S 2.2, pp.~57--60]{Mendelson1997}; the two lemmas below are the
counterparts, in the handbook's semantics, of his Lemma~2 with properties (X)
and (XI) of \S 2.2, and of his Proposition~2.2
\cite[\S 2.2, p.~63; \S 2.4, p.~72]{Mendelson1997}.

\begin{hblemma}{thm:p02-w02c-p02-017-free-for}{proved=annex:p02-w02c-p02-017-free-for}{Substitution of a term free for the variable}
If \(t\) is free for \(x_i\) in \(\mathscr B(x_i)\), then for every structure
\(\mathcal M\) and assignment \(s\),
\[
 \mathcal M,s\entails\mathscr B(t)
 \quad\Longleftrightarrow\quad
 \mathcal M,s[x_i\mapsto\sem{t}^{\mathcal M}_{s}]
   \entails\mathscr B(x_i).
\]
Here \(\mathscr B(t)\) is the replacement of
\cref{def:p02-w02c-p02-017-language}, not the capture-avoiding operation of
\cref{def:p02-fol-capture-avoiding-substitution}.
\end{hblemma}

\begin{hblemma}{thm:p02-w02c-p02-017-soundness}{proved=annex:p02-w02c-p02-017-soundness}{Theorems of a first-order theory hold in its models}
Let \(K\) be a first-order theory and \(\mathcal M\) a structure in which
every proper axiom of \(K\) is satisfied by every assignment.  Then every
theorem of \(K\) is satisfied in \(\mathcal M\) by every assignment.  In
particular, every theorem of a first-order predicate calculus is satisfied in
every structure by every assignment.
\end{hblemma}

The lemma is about theorems, that is, deductions from the empty set.  It does
not extend to \(\Gamma\vdash_K\mathscr C\) read as assignment-wise
consequence: in the countermodel of the proof of the next theorem, the
deduction \(A^1_1(x_1)\vdash_K(\forall x_1)A^1_1(x_1)\) leads from a
hypothesis that an assignment satisfies to a conclusion that the same
assignment does not satisfy.

\begin{hbtheorem}{thm:p02-w02c-p02-017-failure}{proved=annex:p02-w02c-p02-017-failure}{The unmodified deduction theorem fails}
Let \(K\) be a first-order predicate calculus whose language contains the
predicate letter \(A^1_1\).  Then
\(A^1_1(x_1)\vdash_K(\forall x_1)A^1_1(x_1)\), but
\(A^1_1(x_1)\Rightarrow(\forall x_1)A^1_1(x_1)\) is not a theorem of \(K\).
Every deduction of \((\forall x_1)A^1_1(x_1)\) from \(\{A^1_1(x_1)\}\) violates
the hypothesis of \cref{thm:p02-w02c-p02-017-prop25} for the discharged wf
\(A^1_1(x_1)\).
\end{hbtheorem}

Mendelson gives this example with a countermodel of the same kind
\cite[\S 2.4, p.~73]{Mendelson1997}.  The theorem shows that the hypothesis of
Proposition~2.5 cannot simply be dropped; it does not show that the hypothesis
is necessary in every instance, and
\cref{thm:p02-w02c-p02-017-cor27} gives a case, a closed discharged wf, in
which it holds automatically.

\subsection{Relationships and sources}
\label{sec:p02-w02c-p02-017-relations}

\emph{The registered sheet.}  The calculus on
\texttt{sheet:\allowbreak hilbert-fol}, with the derivations
of \cref{def:p02-hilbert-fol-dependencies}, differs from Mendelson's theories
in four respects.  First, its generalization rule may be applied only when the
variable is free in no premise on which the line depends, so the restriction
is part of its derivability relation and its deduction theorem
\cref{thm:p02-hilbert-fol-deduction} has no extra hypothesis; in Mendelson's
theories Gen is unrestricted, \(\mathscr B\vdash_K(\forall x_i)\mathscr B\)
holds for every \(\mathscr B\), and the restriction appears only as the
hypothesis of Proposition~2.5.  Second, the sheet's instantiation schema uses
the total capture-avoiding substitution of
\cref{def:p02-fol-capture-avoiding-substitution}, which renames bound
variables, whereas (A4) uses replacement of occurrences together with the
side condition that \(t\) be free for \(x_i\).  Third, the sheet takes
\(\exists\) as primitive with its own schema and rule, while Mendelson defines
\((\exists x)\mathscr B\) as \(\neg(\forall x)\neg\mathscr B\).  Fourth,
Mendelson's theories may have proper axioms; the sheet has none.  The
universal schemas of the sheet correspond to (A1)--(A5); its dependency-
sensitive generalization rule is not Mendelson's.  Results proved for one
system pass to the other only through explicit derivations.

\emph{Kleene's predicate calculus.}  Kleene, like Mendelson, puts no
condition on assumption formulas in his quantifier rules and states the
restriction as the hypothesis of his deduction theorem, that free variables
be held constant for the discharged assumption
(\cref{def:p02-w02-p02-025-kleene-dependence,thm:p02-w02-p02-025-kleene-deduction}).
His rules have the forms ``from \(C\supset A(x)\) infer
\(C\supset\forall xA(x)\)'' and the existential rule 12, so his ``varied''
variables are those of applications of these two rules, while Mendelson's
condition concerns Gen alone.  Both dependence relations are defined by the
same two clauses, an assumption step and inheritance through rule
applications.

\emph{The propositional system.}  (A1)--(A3) with MP are the axioms and rule
of Mendelson's propositional system \(L\)
\cite[\S 1.4, p.~35]{Mendelson1997}, whose deduction theorem
(Proposition~1.9) has no variable condition
\cite[\S 1.4, p.~37]{Mendelson1997}; the handbook's propositional Hilbert
calculus and its deduction theorem are \cref{thm:p02-hilbert-deduction}.  In
a first-order theory the propositional proof of
\cref{thm:p02-w02c-p02-017-prop25} goes through unchanged for Hyp, axiom and
MP steps; only the Gen case needs the hypothesis, (A5), and
Proposition~2.4.

\emph{Sources.}  Mendelson, fourth edition: formal theories, deductions and
consequence, \S 1.4, pp.~34--35; the system \(L\), p.~35, and Lemma~1.8,
p.~36; terms, wfs and the defined connectives and \(\exists\), \S 2.1,
p.~52; free and bound occurrences, p.~53; the notation \(\mathscr B(t)\) and
the free-for relation, p.~54; first-order languages, \S 2.2, p.~57; closed
wfs, p.~58; first-order theories, (A1)--(A5), MP and Gen, \S 2.3, pp.~69--70,
and the examples for the side conditions, pp.~70--71; failure of the
unmodified deduction theorem, dependence and its example, \S 2.4, p.~73;
Propositions 2.4 and 2.5, p.~74; Corollaries 2.6 and 2.7 and the extension
of Propositions 2.4--2.7, p.~75 \cite{Mendelson1997}.  Mendelson states
Corollaries 2.6 and 2.7 without proof and leaves the extension as an
exercise; the proofs in the annex are written for this section.
